\subsection{Equations of Lines and Planes}


\subsubsection{Lines}
\content{
\vspace{3in}
}{% presentation
}{% nopresentation
}{% comments
Demonstrate that a a line is determined by an initial point, and a direction
vector, and that this gives rise to the equation 
$$ \vec{r}=\vec{r}_0 + t\vec{v}. $$
}

\content{
Parametric equations of a line, where $(x_0,y_0,z_0)$ is a point on the line,
and $\la a,b,c\ra$ is a vector parallel to the line.  
$$ x = x_0+at\quad y=y_0+bt\quad z=z_0+ct. $$
}{% presentation
\vspace{1in}
}{% nopresentation
\vspace{1in}
}{% comments
Show that this can be derived quickly from the vector equation above. 
}

\content{
\begin{example} Find a vector equation and the parametric equations of the line
passing through $(1,3,2)$ and parallel to $\la 1,4,-2\ra$.
\end{example}
}{% presentation
\vspace{2in}
}{% nopresentation
\vspace{1in}
}{% comments
}

\content{
Symmetric equations of a line where $(x_0,y_0,z_0)$ is a point on the line, and
$\la a,b,c\ra$ is a vector parallel to the line. 
$$ \frac{x-x_0}{a} = \frac{y-y_0}{b} = \frac{z-z_0}{c}. $$
}{% presentation
\vspace{1in}
}{% nopresentation
\vspace{1in}
}{% comments
Notice that is one of $a$, $b$, or $c$ is zero, then these symmetric equations
look a little different. 
}

\content{
\begin{example} Find the symmetric equations of the line whose vector equation
is given by 
$$ \vec{r}(t) = (3,2,1) + t\la -1,0,3\ra. $$
\end{example}
}{% presentation
\vspace{2in}
}{% nopresentation
\vspace{1in}
}{% comments
}

\content{
\begin{example} Show that the lines $L_1$ and $L_2$ whose equations are 
\begin{eqnarray*} L_1:  x =1+t,\quad y= -2+3t,\quad z = 4-t \\
L_2: x=2s,\quad y = 3+s,\quad z=-3+4s
\end{eqnarray*}
Are skew lines (that is, they are not parallel, and do not intersect).
\end{example}
}{% presentation
\vspace{2in}
}{% nopresentation
\vspace{2in}
}{% comments
}

\content{
\begin{example} Line segments.
\end{example}
}{% presentation
\vspace{2in}
}{% nopresentation
\vspace{2in}
}{% comments
Derive the formula for a line segment, by drawing two vectors $\vec{a}$ and
$\vec{b}$, and demonstrating that all the points between them 
can be expressed through the equation 
$$ \vec{r}(t) = (1-t)\vec{a}+t\vec{b} $$
}


\subsubsection{Planes}

\content{
\vspace{3in}
}{% presentation
}{% nopresentation
}{% comments
Show that a plane going through a point $\vec{r}_0$ is the set of all vectors
$\vec{r}$ such that $\vec{r}-\vec{r}_0$ is perpendicular to the normal vector
$\vec{n}$. That is 
$$ a(x-x_0)+b(y-y_0)+c(z-z_0)=0.$$
}

\content{
\begin{example} Find the equation of the plane which passes through the point
$(2,4,-1)$ with normal vector $\vec{n}=\la 2,3,4\ra $.
\end{example}
}{% presentation
\vspace{2in}
}{% nopresentation
\vspace{2in}
}{% comments
To sketch this plane, find the intercepts at the axes and draw a triangle.
}

\content{
\begin{example} Find the equation of the plane which passes through the points 
$P(1,3,2)$, $Q(3,-1,6)$ and $R(5,2,0)$.
\end{example}
}{% presentation
\vspace{3in}
}{% nopresentation
\vspace{2in}
}{% comments
}


\content{
\begin{example}
Find the point at which the line $x = 2+3t$, $y=-4t$, $z=5+t$ intercepts the
plane $4x+5y-2z=18$.
\end{example} 
}{% presentation
\vspace{2in}
}{% nopresentation
\vspace{1in}
}{% comments
Simply insert the line equations into plane equation and solve for $t$.
}

\content{
\begin{example} Find the angle between the two planes $x+y+z=1$ and $x-2y+3z=1$. 
\end{example}
}{% presentation
\vspace{2in}
}{% nopresentation
\vspace{2in}
}{% comments
Note that the angle between two planes is defined to be the angle between their
normal vectors. 

Answer is $\cos^{-1}\left(\frac{2}{\sqrt{42}}\right) \approx 72^\circ$.
}

\subsubsection{Distances}

\content{
\vspace{3in}
}{% presentation
}{% nopresentation
}{% comments
Show that the distance from a point $P$ to a plane $D$ is given by taking any
point $P_0$ in the plane, and computing the length of the projection of the
vector $\vec{P_0P}$ onto the normal vector of $D$.
}

\content{
\begin{example} Find the distance from the point $(4,0,1)$ to the plane
$x-y+z=1$.
\end{example}
}{% presentation
\vspace{3in}
}{% nopresentation
\vspace{2in}
}{% comments
}

\content{
\begin{example} Find the distance between the two planes $10x+2y-2z=5$ and 
$5x+y-z=1$. 
\end{example}
}{% presentation
\vspace{3in}
}{% nopresentation
\vspace{2in}
}{% comments
Notice that these two planes are parallel. So we can take any point on one plane
and compute it's distance to the other plane. 
}
